\chapter{The machine-learning toolbox}
\label{ch:toolbox}

\begin{tldr}
\begin{itemize}
\item A model is a function with millions of adjustable numbers (\textbf{weights}). Training nudges them until the
      outputs match the labels.
\item Image models are \textbf{CNNs}: a pretrained \textbf{backbone} turns an image into a feature vector, a small
      \textbf{head} turns features into the 12 outputs.
\item Knee MRI is fed slice by slice (\textbf{2D}) or three neighbouring slices at a time (\textbf{2.5D}); the
      slice results are \textbf{pooled} into one answer per study.
\item Honest testing needs data the model never trained on: \textbf{cross-validation}. Averaging several models
      (\textbf{ensembling}) almost always helps a little.
\end{itemize}
\end{tldr}

\section{A model is a function with knobs}

\begin{figure}[H]
\centering
\begin{tikzpicture}[x=0.86mm,y=1mm]
  \node[draw=grid, fill=white, rounded corners=1.5mm, minimum width=24mm, minimum height=16mm, align=center,
        font=\footnotesize] (x) at (0,0) {\faImage\\ images\\ of one study};
  \node[draw=accent, fill=keyfill, rounded corners=2mm, minimum width=44mm, minimum height=22mm] (f) at (45,0) {};
  \node[font=\small\bfseries] at (45,6) {$f_{\theta}$};
  \foreach \i in {0,...,6} {
    \draw[accent, line width=0.6pt] ({29+\i*5.3},-5) circle (1.6);
    \draw[accent, line width=0.8pt] ({29+\i*5.3},-5) -- ++({90-\i*37}:1.6);
  }
  \node[font=\scriptsize, text=muted] at (45,-10) {millions of knobs $\theta$ (weights)};
  \node[draw=grid, fill=white, rounded corners=1.5mm, minimum width=24mm, minimum height=16mm, align=center,
        font=\footnotesize] (y) at (90,0) {12 numbers\\ (one per finding)};
  \draw[-{Stealth[length=2mm]}, line width=0.8pt, muted] (x) -- (f);
  \draw[-{Stealth[length=2mm]}, line width=0.8pt, muted] (f) -- (y);
  \node[draw=them, fill=whyfill, rounded corners=1.5mm, font=\footnotesize, align=center] (l) at (90,-22) {labels\\ (what we want)};
  \draw[{Stealth[length=2mm]}-{Stealth[length=2mm]}, them, line width=0.8pt, dashed] (y) -- (l)
     node[midway, right, font=\scriptsize] {compare};
  \draw[-{Stealth[length=2mm]}, them, line width=0.8pt] (l.west) -- (60,-22) -- (52,-13)
     node[pos=0.3, below, font=\scriptsize, text=them!85!black] {turn the knobs a little};
\end{tikzpicture}
\caption{Training in one picture: compare the outputs with the labels, then turn every knob a little in the
direction that reduces the difference. Repeat many thousands of times.}
\end{figure}

\begin{wordcheck}
\term{Weights} (parameters): the adjustable numbers inside the model. EfficientNet-B0, our main model, has about
5 million; the winners' models have 30--120 million. \term{Training}: adjusting them. \term{Inference}: using the
trained model on new data.
\end{wordcheck}

\section{How a CNN sees: filters}

A \term{convolutional neural network} (CNN) looks at an image through many small \term{filters}. A filter is a
3\,$\times$\,3 grid of weights that slides over the image; at each position it multiplies and adds, producing one
number. The result is a new image, a \term{feature map}, which is bright where the filter's pattern occurs.

\begin{figure}[H]
\centering
\begin{tikzpicture}[x=4.2mm,y=4.2mm]
  % input image 7x7
  \foreach \i in {0,...,6} \foreach \j in {0,...,6} {
    \pgfmathsetmacro{\v}{(\i<3) ? 15 : 75}
    \fill[black!\v] (\i,\j) rectangle ++(1,1);
  }
  \draw[grid] (0,0) grid (7,7);
  \draw[them, line width=1.4pt] (2,3) rectangle (5,6);
  \node[font=\footnotesize, align=center] at (3.5,-1.4) {image (dark | bright edge)};
  % filter
  \begin{scope}[xshift=10.5cm/4.2]
  \end{scope}
  \foreach \i/\w in {0/-1,1/0,2/1} \foreach \j in {0,1,2} {
    \fill[white] (\i+9.5,\j+3) rectangle ++(1,1);
    \node[font=\scriptsize] at (\i+10,\j+3.5) {\w};
  }
  \draw[them, line width=1.4pt] (9.5,3) rectangle (12.5,6);
  \draw[grid] (9.5,3) grid (12.5,6);
  \node[font=\footnotesize, align=center] at (11,1.6) {filter\\ ``vertical edge''};
  \node[font=\Large] at (8.25,4.5) {$\ast$};
  \node[font=\Large] at (13.75,4.5) {$=$};
  % feature map 5x5
  \foreach \i in {0,...,4} \foreach \j in {0,...,4} {
    \pgfmathsetmacro{\v}{(\i==1 || \i==2) ? 85 : 5}
    \fill[us!\v] (\i+15,\j+1) rectangle ++(1,1);
  }
  \draw[grid] (15,1) grid (20,6);
  \node[font=\footnotesize, align=center] at (17.5,-1.4) {feature map\\ (bright = edge found)};
\end{tikzpicture}
\caption{One filter finds vertical edges. A CNN learns hundreds of filters per layer and stacks dozens of layers:
early layers find edges, middle layers textures, deep layers whole structures (``a meniscus'', ``fluid'').}
\end{figure}

\section{Backbone and head}

\begin{figure}[H]
\centering
\begin{tikzpicture}[x=0.84mm,y=1mm]
  \fill[dark!85] (0,-12) rectangle (24,12);
  \fill[white!65!dark] (12,0) circle (6);
  \node[font=\footnotesize] at (12,-16) {one input image};
  % backbone: shrinking blocks
  \foreach \i/\h in {0/22,1/18,2/14,3/10,4/7} {
    \fill[us!\the\numexpr 25+\i*15\relax, rounded corners=0.6mm] ({32+\i*9},{-\h/2}) rectangle ++(6,\h);
  }
  \draw[decorate, decoration={brace, amplitude=2mm}, line width=0.6pt] (32,14) -- (74,14)
     node[midway, above=2.5mm, font=\footnotesize] {\textbf{backbone} (pretrained CNN)};
  \node[font=\scriptsize, text=muted, align=center] at (53,-16) {edges $\rightarrow$ textures $\rightarrow$ structures};
  % feature vector
  \foreach \k in {0,...,9} \fill[us!80] (80,{-9+\k*1.9}) rectangle ++(3,1.5);
  \node[font=\scriptsize, align=center] at (81.5,-16) {feature\\ vector};
  % head
  \node[draw=them, fill=whyfill, rounded corners=1.5mm, minimum width=14mm, minimum height=12mm, font=\footnotesize,
        align=center] (h) at (99,0) {head};
  \draw[-{Stealth}, muted] (84,0) -- (h);
  \foreach \k in {0,...,11} \fill[them!80] (114,{-11+\k*1.9}) rectangle ++(3,1.5);
  \draw[-{Stealth}, muted] (h) -- (113,0);
  \node[font=\scriptsize, align=center] at (115.5,-16) {12 outputs};
  \draw[-{Stealth}, muted] (25,0) -- (31,0);
\end{tikzpicture}
\caption{Every image model in this script has this shape. The backbone is big and pretrained; the head is small
and new. ``Feature vector'' = a list of numbers that summarises what the backbone saw (1,280 numbers for
EfficientNet-B0).}
\label{fig:backbone}
\end{figure}

\subsection*{Pretrained backbones}

Training a backbone from scratch needs millions of images. Instead everyone starts from a backbone trained on
\term{ImageNet} (1.2 million photos of everyday things in 1,000 classes). Its early filters (edges, textures)
are useful for any image, including MRI. We then \term{fine-tune}: keep training all weights on our task with a
small learning rate. This is called \term{transfer learning}.

\begin{center}\small
\begin{tabularx}{\linewidth}{lrl>{\raggedright\arraybackslash}X}
\toprule
Backbone & Weights & Used by & Speed on a T4 GPU (ours, 256\,px) \\
\midrule
EfficientNet-B0 & 5\,M & us & 0.5\,s per batch of 8 studies \\
ConvNeXt-Nano & 15\,M & us & about 1.4$\times$ slower \\
ConvNeXt-T(iny) & 28\,M & winners & 2.5$\times$ slower, needed tricks to fit \\
MaxViT-V2-B, CoAtNet-2 & 100\,M+ & winners, public & far too slow for us on T4s \\
\bottomrule
\end{tabularx}
\end{center}

\section{2D, 2.5D or 3D?}

A knee series is a 3D stack, but the best pretrained backbones are 2D and expect 3 colour channels (red, green,
blue). The common trick:

\begin{figure}[H]
\centering
\begin{tikzpicture}[x=0.86mm,y=1mm]
  \foreach \col/\title/\txt in {0/2D/{one slice,\\ copied into 3 channels}, 50/2.5D/{3 neighbouring slices\\ as the 3 channels}, 100/3D/{the whole stack at once}} {
    \node[font=\small\bfseries] at (\col+12,29) {\title};
    \node[font=\scriptsize, align=center] at (\col+12,-12) {\txt};
  }
  % 2D
  \fill[dark!85, draw=white] (6,0) rectangle ++(14,14);
  % 2.5D
  \foreach \k/\c in {0/red,1/green!70!black,2/blue} {
    \fill[dark!85, draw=\c, line width=1pt] ({52+\k*3},{\k*2.5}) rectangle ++(14,14);
  }
  \node[font=\scriptsize, text=muted] at (62,-4) {R \quad G \quad B};
  % 3D
  \foreach \k in {0,...,9} \fill[dark!85, draw=white, line width=0.3pt] ({102+\k*1.2},{\k*1}) rectangle ++(14,14);
  % verdicts
  \node[font=\scriptsize, text=wordline, align=center] at (12,-20) {\cmark\ simple, pretrained\\ \xmark\ no depth context};
  \node[font=\scriptsize, text=wordline, align=center] at (62,-20) {\cmark\ pretrained\\ \cmark\ a little depth};
  \node[font=\scriptsize, text=wordline, align=center] at (112,-20) {\cmark\ full depth\\ \xmark\ few good pretrained models};
\end{tikzpicture}
\caption{Three ways to feed slices. We used 2.5D (8 slice triplets per series). The winners found that once
the model knows where to look, plain 2D works as well (Chapter~\ref{ch:winmodels}).}
\label{fig:25d}
\end{figure}

\section{From many slices to one answer: pooling}

A study gives many images (we used 8 triplets $\times$ 3 series = 24). Each goes through the backbone and becomes a
feature vector. They must be combined into \emph{one} prediction per study: this is \term{pooling}.

\begin{figure}[H]
\centering
\begin{tikzpicture}[x=1mm,y=1mm]
  \foreach \k/\w in {0/0.03,1/0.04,2/0.05,3/0.40,4/0.35,5/0.06,6/0.04,7/0.03} {
    \fill[dark!85] ({\k*10.5},0) rectangle ++(8.5,8.5);
    \fill[us!90] ({\k*10.5},14) rectangle ++(8.5,{\w*30});
    \node[font=\scriptsize] at ({\k*10.5+4.25},{16+\w*30}) {\w};
  }
  \fill[them] (35.5,2.5) circle (1.3); \fill[them] (46.5,3.5) circle (1.1);
  \node[font=\footnotesize, anchor=east] at (-2,4.5) {slices};
  \node[font=\footnotesize, anchor=east] at (-2,18) {weights};
  \node[font=\scriptsize, align=left, anchor=west] at (89,12) {\textbf{attention pooling}: a tiny network gives\\
     each slice a weight (they sum to 1);\\ the study feature is the weighted average.\\[1mm]
     Here the two slices that show the tear\\ (orange dots) get 75\,\% of the weight.};
\end{tikzpicture}
\caption{A tear is visible on two or three slices out of 24. A plain average would dilute it; attention pooling lets
the model focus on the slices that matter. Because only the study has a label, not each slice, this setting is
called \term{multiple-instance learning}.}
\label{fig:attnpool}
\end{figure}

\section{From numbers to probabilities}

The head outputs 12 raw numbers called \term{logits} (any value from $-\infty$ to $+\infty$). Each one is squeezed
into a probability by the \term{sigmoid} function $\sigma(z) = 1/(1+e^{-z})$:

\begin{figure}[H]
\centering
\begin{tikzpicture}
\begin{axis}[width=8cm, height=4.6cm, xmin=-6, xmax=6, ymin=0, ymax=1, xlabel={logit $z$}, ylabel={$\sigma(z)$},
    ymajorgrids, xmajorgrids, xtick={-6,-4,...,6}, ytick={0,0.5,1}, xlabel style={font=\footnotesize},
    ylabel style={font=\footnotesize}]
  \addplot[us, line width=1.4pt, domain=-6:6, samples=100] {1/(1+exp(-x))};
  \addplot[only marks, mark=*, mark size=2pt, them] coordinates {(2.6,0.931)};
  \node[font=\scriptsize, anchor=west, align=left] at (axis cs:-5.7,0.8) {Study K, ACL:\\ $z=2.6 \Rightarrow p=0.93$};
\end{axis}
\end{tikzpicture}
\caption{The sigmoid. Each of the 12 findings has its own sigmoid, so all 12 can be high at the same time (a
\emph{multi-label} model).}
\end{figure}

\section{The loss: how wrong is the model?}

Training needs one number that says how bad the current predictions are: the \term{loss}. For yes/no labels the
standard is \term{binary cross-entropy} (BCE):
\[
\text{BCE}(p, y) = -\big[\, y\log p + (1-y)\log(1-p) \,\big]
\]
where $p$ is the predicted probability and $y$ the label.

\begin{figure}[H]
\centering
\begin{tikzpicture}
\begin{axis}[width=9cm, height=5cm, xmin=0.01, xmax=0.99, ymin=0, ymax=3.2, xlabel={predicted probability $p$},
    ylabel={loss}, ymajorgrids, xlabel style={font=\footnotesize}, ylabel style={font=\footnotesize},
    legend style={at={(0.5,0.98)}, anchor=north}]
  \addplot[them, line width=1.3pt, domain=0.02:0.98, samples=100] {-ln(x)};
  \addplot[us, line width=1.3pt, domain=0.02:0.98, samples=100] {-ln(1-x)};
  \addplot[third, line width=1.3pt, dashed, domain=0.02:0.98, samples=100] {-(0.7*ln(x)+0.3*ln(1-x))};
  \addplot[only marks, mark=*, mark size=2pt, third] coordinates {(0.7,0.611)};
  \legend{{label $y=1$}, {label $y=0$}, {soft label $y=0.7$}}
\end{axis}
\end{tikzpicture}
\caption{BCE punishes confident mistakes hard: predicting 0.02 for a true positive costs 3.9. With a
\term{soft label} such as 0.7 (``probably yes''), the loss is smallest exactly at $p=0.7$ (dot). That is how
uncertain labels from reports can be used directly.}
\label{fig:bce}
\end{figure}

\begin{keyidea}
Soft labels let the model learn ``probably yes'' instead of a forced yes or no. Both our pipeline and the
winners' trained on soft labels between 0 and 1.
\end{keyidea}

\begin{wordcheck}
\term{Masked loss}: when a label is unknown (missing), it is simply left out of the loss, so it neither rewards nor
punishes the model. Our training code did this for every missing label.
\end{wordcheck}

\section{The training loop}

\begin{figure}[H]
\centering
\begin{tikzpicture}[x=1mm,y=1mm,
  st/.style={draw=accent, fill=keyfill, rounded corners=2mm, minimum width=30mm, minimum height=11mm,
             align=center, font=\footnotesize},
  ar/.style={-{Stealth[length=2mm]}, line width=0.9pt, accent}]
  \node[st] (a) at (0,18) {\textbf{1. batch}\\ take 8 studies};
  \node[st] (b) at (48,18) {\textbf{2. forward}\\ predict 8 $\times$ 12};
  \node[st] (c) at (48,-6) {\textbf{3. loss}\\ compare with labels};
  \node[st] (d) at (0,-6) {\textbf{4. backward + update}\\ nudge all weights};
  \draw[ar] (a) -- (b); \draw[ar] (b) -- (c); \draw[ar] (c) -- (d); \draw[ar] (d) -- (a);
  \node[font=\footnotesize, align=left, anchor=west] at (72,6) {one turn = one \textbf{step}\\
     all training studies once = one \textbf{epoch}\\[1mm]
     ours: 8 studies per step,\\ \ca 440 steps per epoch, 10 epochs,\\ \ca 33--44 minutes on one T4 GPU};
\end{tikzpicture}
\caption{The loop. ``Backward'' computes, for every weight, which direction lowers the loss (the
\term{gradient}); the \term{optimizer} (we and the winners used \term{AdamW}) takes a small step that way.}
\end{figure}

How big a step? That is the \term{learning rate}. It usually starts small (\term{warm-up}), then decreases
smoothly (\term{cosine decay}):

\begin{figure}[H]
\centering
\begin{tikzpicture}
\begin{axis}[width=9cm, height=4.4cm, xmin=0, xmax=10, ymin=0, ymax=3.4, xlabel={epoch},
    ylabel={learning rate ($\times10^{-4}$)}, ymajorgrids, xlabel style={font=\footnotesize},
    ylabel style={font=\footnotesize}]
  \addplot[us, line width=1.4pt, domain=0:1, samples=20] {3*x};
  \addplot[us, line width=1.4pt, domain=1:10, samples=60] {1.5*(1+cos(deg(pi*(x-1)/9)))};
  \node[font=\scriptsize, anchor=west] at (axis cs:1.5,0.5) {ours: peak $3\cdot10^{-4}$, 1 warm-up epoch, cosine};
\end{axis}
\end{tikzpicture}
\caption{Our learning-rate schedule. The winners used a lower, constant rate for one model and a higher one with
a ``one-cycle'' schedule for another (Chapter~\ref{ch:winmodels}).}
\end{figure}

\section{Overfitting and augmentation}

A model with millions of knobs can simply \emph{memorise} its training studies. It then looks perfect on them and
poor on new ones: \term{overfitting}.

\begin{figure}[H]
\centering
\begin{tikzpicture}
\begin{axis}[width=8.5cm, height=4.6cm, xmin=0, xmax=20, ymin=0, ymax=1, xlabel={epoch}, ylabel={loss},
    ytick=\empty, ymajorgrids, xlabel style={font=\footnotesize}, ylabel style={font=\footnotesize}]
  \addplot[us, line width=1.3pt, domain=0:20, samples=60] {0.85*exp(-x/4)+0.05};
  \addplot[them, line width=1.3pt, domain=0:20, samples=60] {0.75*exp(-x/3.5)+0.3+0.012*x};
  \draw[ink, dashed] (axis cs:8,0) -- (axis cs:8,1);
  \node[font=\scriptsize, anchor=west, text=us] at (axis cs:14,0.13) {training};
  \node[font=\scriptsize, anchor=west, text=them] at (axis cs:14,0.58) {validation};
  \node[font=\scriptsize, anchor=south west, align=left] at (axis cs:8.2,0.75) {stop around here:\\ beyond, it memorises};
\end{axis}
\end{tikzpicture}
\caption{Typical curves. Training loss keeps falling; the loss on unseen studies stops improving and rises.}
\end{figure}

\term{Augmentation} fights overfitting: every time a study is shown, it is changed a little at random, so the
model never sees exactly the same image twice.

\begin{figure}[H]
\centering
\begin{tikzpicture}
  \foreach \x/\lab in {0/original, 3.1/rotated, 6.2/zoomed, 9.3/brighter, 12.4/mirrored} {
    \node[font=\scriptsize] at (\x+1.25,-0.35) {\lab};
    \draw[grid] (\x,0) rectangle ++(2.5,3);
  }
  \begin{scope}[shift={(0.05,0.1)}, scale=0.235] \kneesagittal \end{scope}
  \begin{scope}[shift={(3.35,0.25)}, scale=0.22, rotate=8] \kneesagittal \end{scope}
  \begin{scope}[shift={(6.2,0)}] \clip (0,0) rectangle (2.5,3); \begin{scope}[shift={(-0.35,-0.45)}, scale=0.3] \kneesagittal \end{scope} \end{scope}
  \begin{scope}[shift={(9.35,0.1)}, scale=0.235] \kneesagittal \fill[white, opacity=0.35] (0,0) rectangle (10.5,12); \end{scope}
  \begin{scope}[shift={(14.85,0.1)}, xscale=-0.235, yscale=0.235] \kneesagittal \end{scope}
\end{tikzpicture}
\caption{Augmentations we used: small rotation, zoom and shift, brightness/contrast/gamma changes, and mirroring
the whole knee left--right. All are applied on the GPU, per series, at random.}
\label{fig:augment}
\end{figure}

\section{Validation: testing honestly}

To know whether a change helps, you need studies the model did not train on. \term{K-fold cross-validation}
splits the training studies into $k$ parts (\term{folds}), trains $k$ models, each with one fold held out:

\begin{figure}[H]
\centering
\begin{tikzpicture}[x=0.88mm,y=1mm]
  \foreach \m in {0,...,4} {
    \node[font=\scriptsize, anchor=east] at (-2,{-\m*6}) {model \m};
    \foreach \f in {0,...,4} {
      \ifnum\f=\m
        \fill[them] ({\f*21},{-\m*6-2.2}) rectangle ++(20,4.4);
        \node[font=\scriptsize, text=white] at ({\f*21+10},{-\m*6}) {validate};
      \else
        \fill[us!70] ({\f*21},{-\m*6-2.2}) rectangle ++(20,4.4);
        \node[font=\scriptsize, text=white] at ({\f*21+10},{-\m*6}) {train};
      \fi
    }
  }
  \node[font=\footnotesize, align=left, anchor=west] at (110,-12) {every study is validated\\ exactly once\\
     $\Rightarrow$ \textbf{out-of-fold (OOF)}\\ \hspace*{4mm}predictions for all 4,349};
\end{tikzpicture}
\caption{5-fold cross-validation. At submission time the 5 models are averaged (a free ensemble).}
\label{fig:kfold}
\end{figure}

\begin{watchout}
If studies from the same hospital are in both the training and the validation folds, the model can recognise the
hospital (its scanner, its image style) and the score looks better than it will be on new hospitals. We split
\textbf{by site}: all studies of one site proxy (report language + scanner vendor, 31 groups) go into the same
fold. This is called \term{GroupKFold}.
\end{watchout}

\section{Ensembles and test-time augmentation}

\begin{itemize}
\item \term{Ensemble}: average the predictions of several models. Their mistakes partly cancel. It helps most when
      the models are \emph{different} (other labels, other inputs), and very little when they are near-copies.
\item \term{Rank averaging}: AUC only cares about order, so instead of averaging probabilities you can convert each
      model's predictions to ranks (1st, 2nd, \ldots) and average those. This stops one over-confident model from
      dominating.
\item \term{Test-time augmentation} (TTA): predict on the original and on a mirrored copy, then average. We used
      mirror TTA.
\end{itemize}

\begin{innumbers}
Our EfficientNet-B0 and ConvNeXt-Nano models (one fold each, same images and labels) scored 0.886 and 0.898 on
the 58 gold studies; their rank average scored 0.903. A small but typical ensemble gain.
\end{innumbers}

\begin{checkpoint}
\begin{itemize}
\done I can draw backbone $\rightarrow$ feature vector $\rightarrow$ head and say which part is pretrained.
\done I know what 2.5D means, what pooling does, and why soft labels work with BCE.
\done I can explain k-fold and why we grouped folds by site.
\quiz Q5.1: Why use 3 neighbouring slices as the 3 channels instead of one slice copied 3 times?
\quiz Q5.2: A label is 0.5 (``don't know''). What prediction minimises its BCE, and what does the model learn from it?
\quiz Q5.3: Two models trained on the same images with the same labels and different random seeds. Will
      ensembling them help a lot? Why?
\end{itemize}
\end{checkpoint}
